\documentclass[11pt,a4paper]{article}
\usepackage[margin=1in]{geometry}
\usepackage{parskip}
\usepackage{amsmath,amssymb,amsthm}
\usepackage{mathtools}
\usepackage{xcolor}
\theoremstyle{definition}
\newtheorem{definition}{Definition}
\newtheorem{example}{Example}
\theoremstyle{plain}
\newtheorem{theorem}{Theorem}
\newtheorem{proposition}{Proposition}
\newtheorem{lemma}{Lemma}
\renewcommand{\thesection}{7.\arabic{section}}
\renewcommand{\thesubsection}{\thesection.\arabic{subsection}}
\begin{document}
\begin{center}
    {\LARGE \textbf{Mathematics Method for Physics--Lecture 07}} \\[0.6em]
    {\large Applications of Cauchy's Integral Formula}
\end{center}
\vspace{-0.5em}
\noindent\rule{\textwidth}{0.5pt}
\vspace{0.8em}

\section{Existence of Derivatives}

If $f(z)$ is analytic in a region, then its derivatives $f'(z)$, $f''(z)$, … of
all orders exist, and each of them is again analytic in that region.

For a point $z_{0}$ inside a simple closed contour $C$ on which $f$ is analytic,

\[
f^{(n)}(z_{0}) = \frac{n!}{2\mathrm{\pi}\mathrm{i}}\oint_{C}\frac{f(z)}{(z - z_{0})^{n + 1}}dz .
\]

The formula comes from Cauchy's integral formula by differentiating under the
integral sign $n$ times.

\section{Cauchy's Inequality}

Suppose $f(z)$ is analytic on and inside the circle $C$: $\vert{}z - z_{0}\vert{} = R$.

Suppose that

\begin{equation}
\vert{}f(z)\vert{} \le M\quad on\ C
\end{equation}

Then

\begin{equation}
\vert{}f^{(n)}(z_{0})\vert{} \le \ \frac{n!M}{R^{n}},
\end{equation}

where

\begin{equation}
M = \max_{C}\vert{}f(z)\vert{}
\end{equation}

\textit{Proof.} By the generalized Cauchy integral formula:

\[
\vert{}f^{(n)}(z_{0})\vert{} = \left\vert{}\frac{n!}{(2\mathrm{\pi}\mathrm{i})}\oint_{C}\frac{f(z)}{(z - z_{0})^{(n + 1)}}dz\right\vert{}
\]

\[
\le \ \frac{n!}{(2\mathrm{\pi})}\oint_{C}\frac{\vert{}f(z)\vert{}}{\vert{}z - z_{0}\vert{}^{(n + 1)}}\vert{}dz\vert{}
\]

On $C$, $\vert{}z - z_{0}\vert{} = R$ and $\vert{}f(z)\vert{} \le M$.

Taking absolute values in the integral formula, and using $\vert{}\int g(z)dz\vert{} \le \int\vert{}g(z)\vert{}\vert{}dz\vert{}$:

\[
\vert{}f^{(n)}(z_{0})\vert{} \le \frac{n!}{2\mathrm{\pi}} \cdot \frac{M}{R^{(n + 1)}} \cdot \oint_{C}\vert{}dz\vert{}
\]

The contour length is $2\mathrm{\pi}R$, so

\[
= \frac{n!M}{(2\mathrm{\pi}R^{(n + 1)})} \cdot 2\mathrm{\pi}R = \frac{n!M}{R^{n}}
\]

\section{Maximum Modulus Principle}

Suppose $f(z)$ is analytic in the open disk and continuous on its closure, and let $C$ be its boundary circle.

Write

\begin{equation}
M = \max_{\zeta \in C}\vert{}f(\zeta)\vert{},
\end{equation}

so that $\vert{}f(\zeta)\vert{} \le M$ for every $\zeta$ on $C$. The maximum-modulus
principle says that no interior point exceeds this boundary maximum:

\begin{equation}
\max_{z \in \bar{D}}\vert{}f(z)\vert{} = \max_{\zeta \in C}\vert{}f(\zeta)\vert{}.
\end{equation}

\textit{Proof.} By Cauchy’s integral formula:

\begin{equation}
\vert{}f(z)\vert{} = \left\vert{}\left(\frac{1}{2\mathrm{\pi}\mathrm{i}}\right)\oint_{C}\frac{f(\zeta)}{(\zeta - z)}d\zeta\right\vert{}
\end{equation}

\[
\le \left(\frac{1}{2\mathrm{\pi}}\right)\oint_{C}\frac{M}{\vert{}\zeta - z\vert{}}\vert{}d\zeta\vert{}
\]

\[
\le \left(\frac{1}{2\mathrm{\pi}}\right) \cdot\left(\frac{M}{d}\right)\oint_{C}\vert{}d\zeta\vert{}
\]

\[
= \left(\frac{1}{2\mathrm{\pi}}\right) \cdot\left(\frac{M}{d}\right) \cdot L
\]

\[
= \left(\frac{L}{(2\mathrm{\pi}d)}\right)M
\]

Here $d$ is the distance from $z$ to the contour $C$, and $L$ is the length of $C$.

Apply Cauchy’s integral formula to the function $[f(z)]^{n}$:

For a fixed interior point $z$ at distance $d$ from the boundary, $\vert{}\zeta - z\vert{} \ge d$ on $C$, and $\vert{}f(\zeta)\vert{} \le M$ on C:

\begin{equation}
\vert{}[f(z)]^{n}\vert{} = \left\vert{}\left(\frac{1}{2\mathrm{\pi}\mathrm{i}}\right)\frac{[f(\zeta)]^{n}}{(\zeta - z)}d\zeta\right\vert{}
\end{equation}

\[
\le \left(\frac{L}{(2\mathrm{\pi}d)}\right)M^{n}.
\]

Take the nth root:

\begin{equation}
\vert{}f(z)\vert{} \le \sqrt[n]{\frac{L}{2\mathrm{\pi}d}}\,M
\end{equation}

Let $n \to \infty$. Since

\begin{equation}
\sqrt[n]{\left(\frac{L}{(2\mathrm{\pi}d)}\right)} \to 1,
\end{equation}

we obtain

\begin{equation}
\vert{}f(z)\vert{} \le M.
\end{equation}

\section{Liouville's Theorem}

If $f(z)$ is analytic on the whole complex plane (entire), and $\vert{}f(z)\vert{} \le M$, then $f(z)$ is constant.

\textit{Proof.} Take arbitrary distinct points $z_{1}$ and $z_{2}$.

Consider

\begin{equation}
I = \left(\frac{1}{2\mathrm{\pi}\mathrm{i}}\right)\oint_{C}\frac{f(z)}{[(z - z_{1})(z - z_{2})]}dz.
\end{equation}

Choose the contour $C$ to be $\vert{}z\vert{} = R$, where $R > \max(\vert{}z_{1}\vert{}, \vert{}z_{2}\vert{})$, and let $R \to \infty$.

By partial fractions:

\begin{equation}
I = \left(\frac{1}{2\mathrm{\pi}\mathrm{i}}\right)\oint_{C}\frac{f(z)}{[(z - z_{1})(z - z_{2})]}dz
\end{equation}

\[
= \left(\frac{1}{2\mathrm{\pi}\mathrm{i}}\right)\oint_{C_{1}}\frac{[\frac{f(z)}{(z - z_{2})}]}{(z - z_{1})}dz + \left(\frac{1}{2\mathrm{\pi}\mathrm{i}}\right)\oint_{C_{2}}\frac{[\frac{f(z)}{(z - z_{1})}]}{(z - z_{2})}dz
\]

By Cauchy's integral formula:

\begin{equation}
I = \frac{f(z_{1})}{z_{1} - z_{2}} + \frac{f(z_{2})}{z_{2} - z_{1}}
= \frac{f(z_{1}) - f(z_{2})}{z_{1} - z_{2}} ,
\end{equation}

an expression independent of $R$.

On $\vert{}z\vert{} = R$ the distances to the two points satisfy

\begin{equation}
\vert{}z - z_{1}\vert{} \ge R - \vert{}z_{1}\vert{}, \qquad
\vert{}z - z_{2}\vert{} \ge R - \vert{}z_{2}\vert{} .
\end{equation}

With $\vert{}f\vert{} \le M$ on $C$ and contour length $2\mathrm{\pi}R$, the same integral is
bounded by

\begin{equation}
\vert{}I\vert{} \le \frac{1}{2\mathrm{\pi}}\oint_{C}\frac{M}{\vert{}z - z_{1}\vert{}\vert{}z - z_{2}\vert{}}\vert{}dz\vert{}
\le \frac{MR}{(R - \vert{}z_{1}\vert{})(R - \vert{}z_{2}\vert{})},
\end{equation}

and the right-hand side tends to $0$ as $R \to \infty$. Hence $I = 0$ for every
pair $z_{1} \ne z_{2}$, so

\begin{equation}
f(z_{1}) - f(z_{2}) = 0 .
\end{equation}

Since $z_{1}$ and $z_{2}$ are arbitrary, $f$ is constant.

\textit{Alternatively}, differentiating Cauchy's formula for $f'$ and letting
$R \to \infty$ gives the same conclusion directly:

\begin{equation}
\vert{}f'(z)\vert{} = \left\vert{}\left(\frac{1}{2\mathrm{\pi}\mathrm{i}}\right)\oint_{\vert{}\zeta - z\vert{} = R}\frac{f(\zeta)}{(\zeta - z)^{2}}d\zeta\right\vert{}
\le \left(\frac{1}{2\mathrm{\pi}}\right)\cdot\left(\frac{M}{R^{2}}\right)\cdot 2\mathrm{\pi}R = \frac{M}{R} \to 0 ,
\end{equation}

so $f'(z) = 0$ everywhere, and $f$ is constant.

\section{Mean Value Formula}

Suppose $f(z)$ is analytic in the region bounded by $C$: $\vert{}z - z_{0}\vert{} = R$, and also on $C$.

Then

\begin{equation}
f(z_{0}) = \frac{1}{2\mathrm{\pi}}\int_{0}^{2\mathrm{\pi}}f(z_{0} + R\mathrm{e}^{\mathrm{i}\theta})d\theta
\end{equation}

\textit{Proof.} By Cauchy’s integral formula:

\begin{equation}
f(z_{0}) = \left(\frac{1}{2\mathrm{\pi}\mathrm{i}}\right)\oint_{C}\frac{f(\zeta)}{\zeta - z_{0}}d\zeta
\end{equation}

With

\begin{equation}
\zeta = z_{0} + R\mathrm{e}^{\mathrm{i}\theta}, \qquad
d\zeta = \mathrm{i}R\mathrm{e}^{\mathrm{i}\theta} d\theta ,
\end{equation}

the integral becomes

\begin{equation}
f(z_{0}) = \frac{1}{2\mathrm{\pi}\mathrm{i}} \int_{0}^{2\mathrm{\pi}}
\frac{f(z_{0} + R\mathrm{e}^{\mathrm{i}\theta})}{R\mathrm{e}^{\mathrm{i}\theta}}
\cdot \mathrm{i}R\mathrm{e}^{\mathrm{i}\theta} d\theta
= \frac{1}{2\mathrm{\pi}} \int_{0}^{2\mathrm{\pi}} f(z_{0} + R\mathrm{e}^{\mathrm{i}\theta}) d\theta .
\end{equation}

\section{Converse of Cauchy's Theorem}

Suppose $f(z)$ is analytic in $D$. The following statements are equivalent:

1) For every closed curve $C$ in $D$:

\begin{equation}
\oint_{C}f(z)dz = 0.
\end{equation}

2) $f(z)$ has an antiderivative in $D$.

\textit{Proof.} Fix $z_{0} \in D$. For $z \in D$, define

\begin{equation}
F(z) = \int_{z_{0}}^{z}f(\zeta)d\zeta
\end{equation}

The zero-integral condition makes this definition independent of the chosen path: any two paths from $z_{0}$ to $z$ form a closed curve when one is followed in reverse, and the integral around that closed curve is zero.

\begin{equation}
F'(z) = f(z)
\end{equation}

To compute $F'(z)$, choose the short straight segment from $z$ to $z + \Delta z$. Path independence gives:

On $\zeta = z + t\Delta z$, $0 \le t \le 1$:

\begin{equation}
\frac{F(z + \Delta z) - F(z)}{\Delta z} = \int_{0}^{1}f(z + t\Delta z)dt
\end{equation}

\[
F'(z) = f(z)
\]

Conversely, with $F'= f$:

\begin{equation}
\oint_{C}f(z)dz = \oint_{C}F'(z)dz = 0
\end{equation}

\[
\lim_{\Delta z \to 0}\frac{F(z + \Delta z) - F(z)}{\Delta z} = f(z)
\]

\end{document}
